\section{Limitations and Responsible Use}
\label{sec:limitations}

\textbf{Reference design and coverage.} The human reference was selected for calibration and challenge analysis and was repeatedly examined during guideline development. Exclusion from gradients and checkpoint selection does not remove that dependence. Its small, uneven type distribution limits generalization, and keyword-prioritized expansion does not represent the recruitment population. The new final-handbook study supplies independent three-coder evidence, but mean typed span F1 is 0.663 and L has only two spans. Its 50-sentence teacher-pool sample does not establish corpus-wide reliability. Historical coding and assisted review use different conditions. The annotations describe advertised requirements, not individual applicants' abilities.

\textbf{Data access and reproducibility.} Public materials allow inspection of labels, scoring of saved predictions, and comparison of the released historical coding and review layers. They do not include the complete final expanded-pool training inputs; the annotation release preserves analytical fields but omits administrative account metadata. Some collection records, historical implementation details, and label-version compatibility checks are incomplete. These gaps limit full retraining and verification. Data Availability specifies component-level access and reuse terms; applications also require source permissions, privacy safeguards, human oversight, and external validation.

\textbf{Experimental comparisons.} Sentence-level splits, residual JobBERT training/development text matches, and the inherited CRF's incomplete selection history limit claims about independent generalization. The supervision studies differ in what they change, and expanded-pool runs use single seeds. Inference budgets vary, while some proxy-reported model identities remain unverified. External datasets also differ in labels, splits, scorers, and selection procedures. Their bootstrap intervals condition on the trained seeds and do not address advertisement clustering or multiple comparisons. Chinese encoder predictions have been rescored, but their small paired difference does not establish an initialization advantage. These comparisons describe the tested settings rather than a general ranking of models or languages.

\section{Conclusion}
\label{sec:conclusion}

Chinese-SkillSpan specifies how to delimit and type competency mentions in Chinese recruitment text, with explicit treatment of coordination, contextual types, and negative cases. Its contribution is a versioned resource and an inspectable annotation procedure. The new three-coder study measures final-handbook consistency separately from historical coding and assisted review. Its remaining boundary and type disagreements show where adjudication and rule clarification are still needed.

On the development-used human reference, Qwen adaptation improved end-to-end extraction under fixed inference rules. The additional Chinese-supervised encoders provide reproducible alternatives, but the paired interval does not establish an ESCO initialization advantage. Boundary-sensitive scores and the type, length, and empty-sentence diagnostics identify where further assessment is needed. These findings support diagnostic use of the released resource. The new agreement results provide evidence about the annotation process, but do not establish model generalization to new advertisement documents or resolve all historical label-version differences.
